\input{estilo}
% Los dos tiempos de un mismo evento: dónde cae sobre el eje es su processing
% time, cuándo lo procesamos; la hora que trae adentro es su event time, cuándo
% ocurrió.
\begin{document}
\begin{figura}\begin{tikzpicture}[
    dp/.style={etiq, draw, rounded corners=2pt, fill=white, align=center, inner sep=2.5pt}]
  \draw[|->] (0,0) -- (8.4,0) node[etiq, right=3pt] {processing time};
  \foreach \p in {(0.5,1.2), (1.6,0.7), (1.3,2.0), (2.9,1.9), (3.4,0.9), (6.3,1.2), (7.2,2.1)}
    \node[evento] at \p {};
  \node[evento] (b) at (4.6,2.2) {};
  \node[dp, anchor=south west] (t) at ($(b)+(0.15,0.15)$) {“britney”\\8:00:00};
  \draw[dashed, draw=suave, line width=0.5pt] (b) -- (4.6,0);
  \draw (4.6,-0.12) -- (4.6,0.12);
  \node[etiq, below=5pt] (pt) at (4.6,0) {8:00:07};
  \node[etiq, text=suave, below=1pt of pt, align=center] {processing time:\\cuándo lo procesamos};
  \node[etiq, text=suave, align=left, anchor=west] (et) at (7.0,3.4) {event time:\\cuándo ocurrió};
  \draw[->, draw=suave, line width=0.5pt] (et.west) -- ($(t.north east)+(-0.1,-0.05)$);
\end{tikzpicture}\end{figura}
\end{document}
